\section{Introduction}

The same alias, the same prompt text, the same harness, days apart, returned completions one to two orders of magnitude faster than nominally \emph{smaller} models on the same task, and produced valid geometry in 4 of 30 attempts where those smaller models produced it in 30 of 30 and 50 of 60. Nothing in the experimental design predicts either result. Both point to the serving path rather than to anything in the study --- an attribution \S4 weighs but cannot decide --- and the first was visible only because the harness happened to log per-invocation duration.

The weights behind an alias are a promise, not a hash. An alias such as \texttt{opus} or \texttt{haiku} is resolved at request time by infrastructure the experimenter cannot inspect: different weights on different days, the same weights through a different decode path, or a mixture --- none observable from inside the runtime. For most applications this does not matter. For a study whose dependent variable is sensitive to which quantization is served, it matters completely, and the sensitivity is not hypothetical: it was measured directly, on an axis pass/fail metrics cannot see.

This paper makes four contributions.

\textbf{C1 --- the precision cliff, in variation rather than viability.} On a value-sensitive constructive task, served-weight quantization measurably moves outputs, but not along the axis the founding hypothesis predicted. Degrading a quantization ladder moves viability and validity by no amount detectable at $n = 50$ per rung --- non-rejection, not equivalence, with 7B viability \emph{inverting} --- while the 2-bit rung collapses the variation a proposal carries, leaving format and geometry intact. The failure went unreported until coordinates rather than scores were examined. What carries the claim is registered and replicated: the echo bound \texttt{q2\_k >= 60\%}, \texttt{q4\_k\_m <= 35\%}, written into the runner before execution and held on five never-sampled seeds (19/24 against 1/17), plus a registered must-differ probe that returned its registered branch. Two controls in \S3 bound rather than extend the claim: a fixed-parent probe cuts the effect's magnitude while leaving its direction and 2-bit endpoint intact, and on that probe's registered centre-only measure the presented lattice is reproduced at 92--98\% at every rung --- but because the probe's parents share one lattice, it speaks to radius variation only, not to whether a better-served proposer would move to a different construction. Both waves' registered decision rules return unusable verdicts (unclassified on wave 1, FAILED on wave 2). Six later registered waves bound generality (\S6, \S8): a hosted capability ladder does not reproduce the failure mode (0/120, DISSOCIATION), the rung \emph{contrast} transfers to a second geometry while the absolute bands do not, a sixth, screen-selected family clears its upper-rung control floor and reproduces the search-activity collapse under a registered floor-gated secondary (its echo primary underpowered), and the transfers that could not be evaluated --- two tasks and five model families below the required competence --- are reported under their own registered floor labels.

\textbf{C2 --- a forensic case study.} A thirty-invocation arm addressed only by alias, where serving-signature and behavioral evidence together \emph{mildly favour} an unattested serving path without deciding it, and where the favouring evidence --- the latency gap --- is a working-log observation the released artifact cannot check. \S4.1 adds that the favouring is unrecoverable: the alias drifted within six days, so no later call can re-test it.

\textbf{C3 --- a non-identifiability argument, conditioned on the runtime observed.} For this agent runtime, the observable set exposed per invocation contains no variable that attests the serving path. We state that set explicitly, show which hypothesis pairs it separates, and name the experiments that would discriminate the two readings behaviorally --- two runnable inside the runtime, one requiring a pinned endpoint outside it. We ran both runnable ones (\S4.1) under a rule locked beforehand, and the result constrains the argument rather than confirming it: the alias's behaviour and serving signature both moved within six days, unannounced, so the hypothesis pair the discriminator was built to separate turns out to presuppose a stable referent it does not have. The claim is conditioned on one vendor and one runtime (\S8); a runtime exposing a serving-path flag or a weights digest would falsify it for that runtime.

\textbf{C4 --- a repair protocol, including its own failures.} The maximal reproducibility such a study \emph{can} achieve --- prompt hashing before sampling, verbatim raw storage, dated alias maps, deterministic local scoring, hash-locked preregistration --- is specified, implemented, and audited against itself: three of five repairs are fully implemented in the released artifact, and we name the two that are not. The audit extends to the one \emph{instrument} we proposed rather than restated: run against \S3's own rows, on a stack where the serving path is fixed by construction, its firing condition fires anyway, and its obvious repair fails across waves, warm-up and hardware. We withdraw it and report the three mechanisms that broke it.

\section{Background and task}

The benchmark is circle packing: place exactly $N$ non-overlapping circles inside the unit square (or a $1 \times a$ rectangle) to maximize the sum of radii. It is a standard target for LLM-guided search, and its value here is that candidates are \emph{scored by an exact local evaluator} --- no model, no judge, no rubric. A proposal is a list of \texttt{[x, y, r]} triples; validity is a finite set of arithmetic inequalities; the objective is a sum. Every source of run-to-run variance therefore lives on the model side of the interface.

\emph{Companion paper.} Two prior documents are cited throughout: the \emph{companion paper} is the tier-substitution forensic study released alongside this one; the \emph{antecedent study} is the combined technical report (\texttt{precision-cliff-paper-combined.md}) whose executed quantization results \S3 condenses. The companion paper establishes why this task is precision-sensitive rather than merely difficult. Proposals concentrate on constructible templates --- a $k \times k$ grid of radius $1/(2k)$, optionally extended with fillers of radius $(\sqrt{2} - 1)/(2k)$ --- whose values are exact rationals and algebraic numbers. The template reached for is predicted in closed form by $k^*(N) = \lfloor \sqrt{N} + \tfrac{1}{2} \rfloor$, the sign of $N - (k^*)^2$ deciding extension versus truncation; across every held-out cell where that rule and the true optimum disagree, the rival optimum was reached in 2 of 34 valid samples. Because predicted values are exact, a proposal either sits on an anchored construction to within rendering precision or it does not, detectable at \texttt{1e-6} --- three orders of magnitude below the $\sim$\texttt{1e-2} gap between competing constructions. These behaviors are also tier-dependent: one tier truncates templates, another perturbs and mixes radii. A study that cannot attest which weights served it cannot attest which regime it measured.

\emph{The instrument, and why the tolerance is registered.} The companion's evaluator scores every row at \textbf{two} tolerances, primary declared in advance. Proposers print six to eight decimals, so an eight-decimal tangency misses exactness by roughly \texttt{5e-9} and a strict \texttt{1e-9} gate reports a correct construction as invalid on rendering grounds alone. Both tolerances are logged for every row (\texttt{valid}, \texttt{valid\_strict\_1e9} in \texttt{arm\_f\_repro.py}), with \texttt{1e-6} primary because it sits far below the $\sim$\texttt{1e-2} separation between rival constructions. Choosing between them after seeing results would have been the easiest way to fabricate a table, which is why the choice is registered rather than argued, and why \S4 reports validity at both.

Several serving-stack variables sit between ``the model'' as named in a paper and the tokens that arrive: weight quantization and activation precision; fast-mode or speculative decode paths; alias rebinding as new dated ids are released; sampling defaults chosen by the runtime rather than the experimenter; and, in agent runtimes specifically, an inherited system prompt plus user-level instruction files that are not part of the task prompt and are not held fixed across time. Only the last is even nominally under the experimenter's control, and only if they never edit their own configuration.

\section{The precision cliff}

This section condenses the executed quantization results of the antecedent study (\texttt{precision-cliff-paper-combined.md}: protocol \S3.4, results \S5.9, summary in contribution 6, scope caveats \S6). No new \emph{ladder} data is reported here --- every rung figure comes from the antecedent's runs. One dataset appears here that the antecedent does not analyse: the fixed-parent dispersion probe (\S3.5, detail in Appendix B), which controls a confound in the antecedent's own echo measurement.

\subsection{What carries this section}

\textbf{Two registered predictions held, and they are what the claim rests on.} The fresh-seed echo bound --- \texttt{q2\_k >= 60\%}, \texttt{q4\_k\_m <= 35\%}, written into \texttt{kaggle\_precision\_sweep\_14b\_fresh.py}'s header and pushed before that run executed --- was observed at \textbf{79\% (19/24) against 6\% (1/17)} on five seeds never sampled in any prior run. The must-differ decision rule, which specified both branches in advance, returned the copying-is-instruction-insensitive branch at \textbf{5 of 5} valid outputs echoing under an explicit prohibition. Around those sit two descriptive replications, labelled as such wherever used: an independently produced imatrix Q2\_K echoing at 75\% against IQ2\_M's 43\%, and a fixed-parent control preserving direction and the 2-bit endpoint at 92\% against 33\% in the matched band.

\textbf{And here is what failed.} The fresh wave's \emph{designated primary} --- F1, an improvement-count prediction --- was refuted. The dispersion probe's wave-1 decision rule returns \textbf{no category} and wave 2's returns \textbf{FAILED}, its ceiling rung never having run for want of a second GPU. The registered scheme-versus-bit-width control returns its \textbf{inconclusive} branch. The 7B ladder refutes the project's founding hypothesis outright. Everything else --- the search-progress statistic, the decomposition, the horizon model, all twenty-nine $p$-values (accounted in Appendix A.4) --- is descriptive or post-hoc and carries no claim.

\textbf{Every 14B figure in this section is recomputed from the raw candidate ledgers, not copied from the antecedent prose; the 7B figures are not.} \texttt{sec3\_ladder\_repro.py}, released with this paper, replays each \texttt{(quantization, seed)} lineage from the coordinate-logged jsonl and re-derives the 14B viability, validity, echo, per-seed, must-differ and invalid-row figures below, plus five of this section's nine Fisher tails ($p = 1.7\times10^{-7}, 3.4\times10^{-6}, 0.007, 0.017, 0.056$). Ledgers: \texttt{sec3\_\bk{}artifacts/\bk{}**/\bk{}candidates\_\bk{}precision\_\bk{}14b.jsonl} (200 rows), \texttt{...\_fresh.jsonl} (100), \texttt{...\_iq2.jsonl} (100), the two \texttt{mustdiffer\_*.jsonl} files, each with a \texttt{provenance.json} recording SHA-256 and byte length of every served weight file. The remaining four Fisher tails ($p = 5.7\times10^{-10}, 0.001, 0.44, 1.0$) and the 7B paragraph are replayed by \texttt{sec3\_7b\_repro.py}. No figure in \S3 requires a reader to compute a tail unaided. The $p = 0.001$ figure is a \textbf{pooled} contrast --- fifteen upper-rung lineages against Q2\_K's five, not rung-versus-rung; the paired Q4\_K\_M-versus-Q2\_K form is $p = 0.048$, and the script prints both.

\textbf{Inferential status of the $p$-values in this section.} Twenty-nine values appear, in four families that must not be pooled, and \textbf{no claim in this paper rests on any of them}; each value's family and the defect that family carries are accounted in Appendix A.4.

Relatedly, ``viability is flat'' and ``validity does not move'' report \emph{non-rejection}, not demonstrated equivalence: no equivalence test was run and $n = 50$ per rung is not powered for one (the Wilson interval on 7/50 is $[0.070, 0.262]$, compatible with a fourfold ratio to the highest rung).

Two notes for a replicator: the running parent advances only on \emph{strict} score improvement, and reimplementing on ties instead undercounts echo by roughly a quarter at the 2-bit rung; and one figure does not reproduce exactly --- near-copy fraction among invalid IQ2\_M rows is reported as 2/22, our replay of the stated definition returns 1/22, a one-row discrepancy at a definitional boundary that affects no claim.

\subsection{Ladder and protocol}

Proposers: \textbf{Qwen2.5-Coder-7B-Instruct} and \textbf{Qwen2.5-Coder-14B-Instruct}, official Qwen GGUF releases, served locally through llama.cpp on $2 \times$ NVIDIA T4, SHA-256 of each weight file recorded in provenance. Task: $N = 26$ in the unit square; every lineage starts from the same seeded valid parent --- 26 uniform-radius circles on a $6 \times 5$ lattice, summed radii 0.89999 (written 0.900). Sampling: temperature 0.8, top-p 0.95, deterministic per-call seed, $\le$1,200 new tokens in a 4,096-token context, one bare chat completion per proposal --- not agentic, removing the orchestration layer as a confound. Per rung: 5 seeds $\times$ 10 generations (50 loop candidates) plus 6 zero-shot probes. The 7B ladder ran five rungs --- FP16 (16.0 effective bits/weight), Q8\_0 (8.5), Q4\_K\_M (4.85), Q3\_K\_M (3.91), Q2\_K (3.35) --- for 250 loop rows and 30 probes; the 14B ladder ran the four quantized rungs (FP16 omitted: $\sim$29.5 GB does not fit $2 \times$ T4) for 200 loop rows and 24 probes. Every candidate is logged live at evaluation time and scored by the same deterministic evaluator at \texttt{1e-6}. The main ladder is a single quantization family --- llama.cpp K-quants, never mixed; the sole deliberate departure is the IQ2 control below.

Effective bits per weight are llama.cpp's nominal figures and \emph{co-vary with the quantization algorithm} rather than indexing it independently --- everything in this section is a \textbf{Q2\_K claim, not yet a bit-width claim}, a distinction the IQ2 control turns from caveat to measurement.

\subsection{At 7B the founding hypothesis failed}

Viability was flat from FP16 down through Q3\_K\_M, a $4\times$ compression: 7/50, 6/50, 9/50, 7/50, all pairwise Wilson intervals overlapping (FP16 vs Q3\_K\_M Fisher $p = 1.0$). The 2-bit rung \emph{inverted}, at 16/50, against the pooled upper four rungs (16/50 vs 29/200, $p = 0.007$ uncorrected, candidate-level) --- not FP16 alone (16/50 vs 7/50, $p = 0.056$, ns). Both contrasts are post-hoc. The mechanism is arithmetic: viability reduces almost entirely to count accuracy --- of 250 proposals only 45 contained exactly 26 circles and 153 contained 24--25 --- the modal count from FP16 through Q3\_K\_M sits one or two short of the parent shown, while Q2\_K's count distribution is broader (tail to 30--36) and centers on the target. Token truncation is ruled out directly: no 24--25-circle proposal came near the 1,200-token cap (max 687 tokens). The antecedent study names this the \textbf{format lottery}. No probe was valid at any rung (0/30); the canonical anchored value was never emitted in 280 opportunities; the best score in the whole sweep, 1.79998 (written 1.800), sits 29\% below the canonical anchored value, 2.5414. At this scale the proposer sits below the floor at which precision could matter.

\subsection{At 14B, a cliff --- in proposal variation, not viability}

Viability runs 22/50, 22/50, 24/50, 19/50 across Q8\_0/Q4\_K\_M/Q3\_K\_M/Q2\_K, all intervals overlapping, against 87/200 versus 38/200 pooled for the scale contrast (Fisher $p = 1.7\times10^{-7}$; the 38/200 denominator is the 7B ladder at the same four rungs, Q8\_0/Q4\_K\_M/Q3\_K\_M/Q2\_K $= 6+9+7+16$, distinct from the 29/200 over FP16--Q3\_K\_M used above). Validity likewise does not move (18/50, 20/50, 19/50, 18/50); recall stays absent (probes 0/24; best score 1.625, 36\% below the 2.5414 anchor). What moved was the capacity to depart from the parent. \textbf{Coordinate-verified parent echo} ran 2/18, 3/20, 3/19 across the upper three rungs (8/57 pooled, 14\%) against \textbf{17/18 (94\%) at Q2\_K} (Fisher $p = 5.7\times10^{-10}$, superseded in magnitude by the fixed-parent control). Every Q2\_K seed echoes at or above half its valid rows (7/7, 4/4, 4/4, 1/1, 1/2 --- the last is exactly 50\%) against at most 2 echoes per seed at any upper rung; 15 of 17 echoes come from three seeds, the remaining two contribute only 1 and 2 valid rows each, so per-seed denominators are too small to support more than the qualitative point. The seed-level improvement contrast (15/15 seeds improving above the boundary vs 1/5 below, $p = 0.001$) is post-hoc and single-run; the candidate-level echo contrast is descriptive only.

\textbf{Data loss, and a coordinate-verified re-execution.} The original 14B run's hosting session expired before the output directory was retrieved, destroying per-candidate jsonl, probe texts and checkpoints; only the verbatim live console log survived, which never recorded coordinates. The rung was re-executed end to end against the same SHA-256-pinned weights, seeds, prompts and sampling parameters, with three logging-only additions: per-candidate coordinates, redundant console echo, and quantitative predictions embedded before execution --- the runners are released (\texttt{sec3\_artifacts/runners/}). This is a \textbf{re-execution, not a replication}: llama.cpp's seeded sampling is deterministic, so console-logged predictions were guaranteed to hold and carry no evidential weight. The falsifiable predictions were the coordinate-level ones --- whether score-matches were verbatim copies or same-sum rearrangements. Both held (Q2\_K echo $\ge$60\%: observed 94\%; each upper rung $\le$35\%: observed 11--16\%), and the check \emph{refuted} the original score-inferred estimates on the upper rungs, reclassifying 4 of 12 presumed echoes as rearrangements, making the verified cliff sharper (14\%$\to$94\%) than the inferred one (21\%$\to$94\%). A public Kaggle kernel, \texttt{precision-redetermin} (owner handle withheld for double-blind review), re-ran 112 of the 224 generations (Q4\_K\_M and Q2\_K) on the original hardware class, and \textbf{all 112 reproduced their archived per-row digests exactly} (verdict SUPPORTED). Two caveats travel with it: the original llama-cpp-python build is unrecoverable (wheel version-pinned 0.3.34), and two earlier executions of the same kernel drew a P100 and correctly returned VOID\_ARCH for every row instead of comparing across architectures --- cross-hardware nondeterminism is Yuan et al.'s finding (\S7), not a defect.

\textbf{Fresh seeds replicate the echo cliff, and a registered prediction fails.} The protocol was run on five never-before-sampled seeds (2222, 3333, 5555, 7777, 9999) at the two decisive rungs, 100 new coordinate-logged candidates. Echo replicates: \textbf{19/24 (79\%) at Q2\_K against 1/17 (6\%) at Q4\_K\_M} (Fisher $p = 3.4\times10^{-6}$, candidate-level, descriptive; the bound's thresholds were written after the destroyed first run's score-inferred estimates of the same contrast were known, so it is a registered replication target rather than a blind prediction), every fresh Q2\_K seed again majority-echo (4/5, 4/5, 5/6, 2/2, 4/6). \textbf{The wave's designated primary was not the echo bound but the improvement count, and it failed}: the runner header labels F1 ``Improvement (seed-level, primary)''; F2, the echo bound, carries no such label. F1 forecast $\le$2/5 Q2\_K seeds improving past baseline; 3/5 did, against 5/5 at Q4\_K\_M (Fisher $p = 0.44$). The antecedent study demotes the improvement collapse to a fragile correlate: the \textbf{echo rate} is the cliff's replicable signature. The $p = 0.001$ seed-level figure quoted earlier in this section is the superseded one.

\textbf{Mechanism: the cliff survives an explicit prohibition.} A registered must-differ probe held the parent fixed and demanded the output not be identical, $\ge$3 circles changed. At Q2\_K the model returned a coordinate-identical copy in \textbf{5 of 5} valid outputs; at Q4\_K\_M, 1 of 5, with all four non-echo outputs scoring above the parent. Under the pre-registered rule ($\ge$50\% echo $\to$ instruction-insensitive; $\le$20\% $\to$ instruction-sensitive; inclusive both bounds; the rule's subject is Q2\_K alone, Q4\_K\_M's 1/5 reported without a branch) Q2\_K falls in the instruction-insensitive branch at the maximum possible rate, while format adherence stays intact.

\textbf{Scheme-graded, not bit-graded: the IQ2 control.} A registered control ran two independently produced 2-bit-class quantizations from a third-party imatrix-calibrated repository, SHA-256-pinned: \textbf{imatrix Q2\_K} ($\sim$3.4 b/w) and \textbf{IQ2\_M} ($\sim$2.7 b/w --- fewer bits, different scheme family, same calibration). Echo: 24/32 (75\%) for imatrix Q2\_K against 12/28 (43\%) for IQ2\_M ($p = 0.017$), must-differ echoes 6/8 versus 0/2, seeds improving 3/5 versus 5/5. The cliff replicates on an independently produced Q2\_K at essentially the official file's matched-seed rate (75\% vs 79\%). The registered bit-width question returns its \textbf{inconclusive branch} --- 43\% falls between the thresholds ($\ge$60\% generalizes, $\le$35\% attributes to algorithm) --- and neither is claimed. The direction points away from raw bit width: 6\% (Q4\_K\_M) $\to$ 43\% (IQ2\_M, $\sim$2.7 b/w) $\to$ 75--79\% (imatrix and official Q2\_K). Read at face value this says a lower-bit quantization from a different family preserves more proposal variation than a higher-bit K-quant, but the fixed-parent control below withdraws that reading (\S3.5).

\textbf{Closing the selection-effect loophole.} Echo rates are computed among \emph{valid} outputs, and a verbatim copy is valid by construction, so a high echo rate is in principle compatible with a model attempting many mutations that all break. Classifying every non-valid row rules this out: there is no garbage --- every failed 14B output at every rung is a fully parsed, malformed circle list, never truncation or non-coordinate text. At official Q2\_K the \emph{invalid} rows are themselves majority near-copies (18/32 re-execution, 15/26 fresh seeds; typically 26 of 27 circles unchanged). Counted over the whole output distribution, official Q2\_K emits copies or near-copies in roughly 70\% of attempts. The scheme gradient reappears in the failures --- near-copy fraction among invalid rows: 0/33 (fresh Q4\_K\_M) $\to$ 2/22 (IQ2\_M) $\to$ 5/18 (imatrix Q2\_K) $\to$ 15/26 and 18/32 (official Q2\_K). A validity-filter explanation predicts mutation-dominated invalid rows; the invalid rows are copy-dominated.

\subsection{A fixed-parent control narrows the contrast}

The echo rates above are measured against each lineage's \emph{running} parent, and those parents are not distributed alike across rungs: Q2\_K lineages rarely improve, so their parent mostly remains the seeded 0.900 grid, while upper rungs climb away from it. Parent quality is confounded with rung, in the direction that flatters the result.

A fixed-parent dispersion probe removes that confound: it runs the same Qwen2.5-Coder-14B at three of the ladder's four rungs, Q4\_K\_M, Q3\_K\_M and Q2\_K (Q8\_0 was registered but never attempted on either wave --- the 8-bit file does not fit the single P100 allocated; both \texttt{provenance.json} files record \texttt{"q8\_0": \{"skipped": "needs 2 gpus, have 1"\}}), holding the parent constant at each of six preset scores from 0.88 to 1.65, logging coordinates and an echo flag per row across two independently seeded waves (v1: 288 rows, 165 valid; v2: 432 rows on out-of-sample salted seeds). Both waves carry their own preregistration (Appendix B); \textbf{wave 2's own registered verdict is FAILED}, because its addendum requires a ceiling rung (q8\_0) that never landed, so every wave-2 descriptive is exploratory.

The registered centres-only echo measure shows \textbf{no cliff} (92--98\% at every rung, both waves), but all six of the probe's ``parents'' share identical centres, differing only in radii --- the probe shows the model one lattice throughout, so this measures whether that lattice is copied, not whether the model would move to a different construction. On the axis the design \emph{can} measure --- radius variation --- the cliff appears: within the fixed lattice, upper rungs perturb radii in 31--51\% of valid rows with no stable ordering between them (reversing between waves), while at Q2\_K almost every valid row is an outright copy, leaving 2/49 (v1) and 5/84 (v2) that perturb radii. Full decomposition, both waves, and the registered analysis table ($JT = 4460.0$ against null mean 4498.5, $p = 0.686$ on the primary echo measure; NED $p = 0.030$; rarefaction $p = 0.953$; quality fork $p = 0.454$) are in \textbf{Appendix B}.

Restricted to the 0.88--0.90 band the loop actually occupies, the loop-vs-fixed-parent contrast is \textbf{6/18 versus 11/12 --- 33\% against 92\%}, not the loop's 15\% against 94\%. Appendix B works out three consequences. The cliff's direction and 2-bit endpoint survive, but its spread does not: the loop's 14\%-to-94\% figure is partly an artifact of where each rung's running parent sits, and the fresh-seed 79\%-vs-6\% figure and the invalid-row near-copy gradient carry the same confound. The IQ2 scheme-gradient ladder (\S3.4) is weakened, since if Q4\_K\_M's matched rate is 33--52\% rather than 6\%, IQ2\_M's 43\% may not be separated from it at all --- so we withdraw the gradient as evidence about bit width and keep it only as a loop-condition description. And the registered upper-rung bound ``$\le$35\%'' held under the loop condition (11--16\%) but would have \textbf{failed} under fixed parents (52--61\%): it was registered against the loop design and is correctly scored there, but is not a general property of those rungs. The probe is not a substitute for the loop measurement, since a fixed parent removes the lineage history, but it is the better-controlled comparison on the one axis where the loop is confounded.

\subsection{What the loop actually did: search progress at the lineage level}

Every measure above is a property of a single call. This subsection reports what the \emph{loop} did with those calls --- post-hoc, appearing in no preregistration, released as \texttt{sec3\_search\_progress.py}. Define an \textbf{accepted improvement} as a valid output scoring strictly above its lineage's running best. On the registered 14B ladder, across five lineages of ten generations each: Q2\_K takes 0, 0, 0, 0, 1 accepted steps (\textbf{1/50} total; final bests 0.900, 0.900, 0.900, 0.900, 1.625); Q3\_K\_M 4, 3, 4, 2, 2 (\textbf{15/50}); Q4\_K\_M 6, 3, 2, 3, 2 (\textbf{16/50}); Q8\_0 2, 2, 5, 1, 4 (\textbf{14/50}). Four of five Q2\_K lineages take zero steps in ten generations. Because an echo can never be an accepted improvement (\texttt{accepted <= valid - echo} by construction, with gaps 0, 1, 1, 2 across the four rungs), most of this contrast is the echo contrast in complement --- the table's addition is the \emph{unit}: a count of steps over a lineage, the quantity a loop's cost is denominated in.

\textbf{The fresh wave's designated primary was a prediction about this quantity, and it failed.} F1 (``Improvement, seed-level, primary''): $\le$2/5 Q2\_K seeds and $\ge$4/5 Q4\_K\_M seeds improving. Observed 3/5 vs 5/5, $p = 0.44$ --- \textbf{F1 failed}. The echo bound F2 carries no primary label. The finer, post-hoc form of the same rows --- steps rather than whether-improved --- gives 3 versus 14, $p = 0.0079$, and does separate the rungs; that gap is analytic flexibility, not a repair of F1, and no claim rests on it. The most this paper claims for the fresh wave: the echo bound held on never-sampled seeds, and its designated primary did not.

\textbf{Quantization gates the frequency of departure; whether it also degrades departure quality is open.} What the 2-bit rung loses is the occasion to depart: 10 departures in 133 valid fixed-parent outputs against Q4\_K\_M's 74 in 164. Whether those rare departures are also \emph{worse} ones is not settled here --- ten departures against a self-set floor of twenty-five is underpowered, and the wave designed to supply more departures returned UNINFORMATIVE by its own control floor (\S8) --- so the question stays open.

This points at a repair and complicates it: if frequency is gated, forcing departure should recover the search --- but the must-differ probe did exactly that and returned 5/5 echoes anyway. Whatever gates departure at 2 bits was, on this probe's evidence, not reachable by instruction.

\textbf{Replication is partial, and the level of the test decides it.} On an exact lineage-level permutation test (distinct from the dispersion probes' JT tests), the registered ladder gives $p = 0.0008$ against the pooled upper three (15,504 splits) and $p = 0.0079$ against Q4\_K\_M alone (252 splits, the floor of a five-versus-five enumeration). The fresh seeds replicate: 3 steps against 14, $p = 0.0079$ (also floor-bound). \textbf{The IQ2 control does not}: imatrix Q2\_K takes 6 steps against IQ2\_M's 16, exact tail $p = 0.135$. A call-level Fisher test on the same counts returns $p = 0.028$; we report the lineage-level figure instead because calls within a lineage are not exchangeable, which costs us the IQ2 replication but avoids banking an unsupported significance.

\textbf{The outcome metric is blind to all of it.} Final best score after ten generations does not separate the rungs anywhere: registered ladder Q2\_K mean 1.045 vs pooled 1.115 ($p = 0.593$) and vs Q4\_K\_M's 1.054 ($p = 0.937$); fresh seeds Q2\_K \emph{higher}, 1.153 vs 0.999 ($p = 0.421$); IQ2 control 1.087 vs 0.922 ($p = 0.318$). This is non-rejection at five lineages per cell, not equivalence --- the score support is lumpy and discrete (0.900, 0.936, 1.040, 1.300, 1.625), so one lucky jump in a Q2\_K lineage lands on the same rung that four small Q4\_K\_M steps climb to. The finding is the \emph{blindness}: an order-of-magnitude collapse in search activity, 1 step against 14--16, leaves the reported outcome statistically indistinguishable at this horizon.

Whether the gap opens at longer horizons is registered here as a \textbf{prediction, not a result} --- we did not run it, and Appendix C reports the horizon-power exercise we did run, including a withdrawn power figure and the structural reason final best score is the wrong dependent variable at any horizon.

\textbf{At 7B the statistic does not replicate, and reverses.} No coordinates were stored at 7B, but accepted improvements are computable. Step counts: FP16 6/50, Q8\_0 2/50, Q4\_K\_M 3/50, Q3\_K\_M 4/50, \textbf{Q2\_K 7/50} --- the 2-bit rung takes the \emph{most} steps of five, permutation tail $p = 0.116$ against pooled others, $p = 0.238$ against Q4\_K\_M. Final best score does not separate the rungs at this scale either ($p = 0.158$ against pooled others, 0.119 against Q4\_K\_M --- completing the six outcome-level tails, $p$ = 0.12--0.94 across ladders, behind the abstract's statement that final best score separates the rungs nowhere; Appendix C). Validity at 7B runs 4--12 of 50 at every rung, so no rung is searching enough for the statistic to have room to separate; the reversal sits with the viability inversion \S3.3 already reports rather than contradicting it. What this establishes: \textbf{the search-step collapse is a 14B observation, not a property of the quantization ladder as such.}

\textbf{What varied, and what did not.} Across every run in this section the inference stack was \emph{constant within each run} --- same llama.cpp build, sampling parameters, harness, and one fixed hardware allocation ($2 \times$ T4 for all four ladders, a single P100 for both dispersion-probe waves, per each \texttt{provenance.json}'s \texttt{gpus} field). Ladder timings and probe timings are therefore never compared to each other. The only thing that changed \emph{inside} a run was which SHA-256-pinned weight file was loaded --- this section measures the effect of \textbf{served-weight quantization}, not a decode path, kernel, batching regime, or speculative-decoding scheme. That degrading this one serving-stack element selectively collapses proposal variation does not establish that other serving-stack variables do the same; that generalization is the hypothesis \S4 finds an unattestable instance of, not a result established here.

The transfer \S3 needs is narrow: a degraded \emph{served quantization} can leave every pass/fail axis intact while removing constructive competence. Measured at $N = 26$ on locally served open weights, one scale pair, one quantization family plus one control --- not at \S4's cells, not inside the agent runtime (\S8).
